\documentclass[conference]{IEEEtran}
\IEEEoverridecommandlockouts
% The preceding line is only needed to identify funding in the first footnote. If that is unneeded, please comment it out.
\usepackage{cite}
\usepackage{amsmath,amssymb,amsfonts}
\usepackage{algorithmic}
\usepackage{graphicx}
\usepackage{textcomp}
\usepackage{xcolor}
\def\BibTeX{{\rm B\kern-.05em{\sc i\kern-.025em b}\kern-.08em
    T\kern-.1667em\lower.7ex\hbox{E}\kern-.125emX}}
\begin{document}

\title{Reconstruction of missing oceanic environmental field data using small-sample, multivariate methods \\
{\footnotesize \textsuperscript{*}Local Prior and Joint Spatiotemporal Learning}
\thanks{Identify applicable funding agency here. If none, delete this.}
}

\author{\IEEEauthorblockN{1\textsuperscript{st} Given Name Surname}
\IEEEauthorblockA{\textit{dept. name of organization (of Aff.)} \\
\textit{name of organization (of Aff.)}\\
City, Country \\
email address or ORCID}
\and
\IEEEauthorblockN{2\textsuperscript{nd} Given Name Surname}
\IEEEauthorblockA{\textit{dept. name of organization (of Aff.)} \\
\textit{name of organization (of Aff.)}\\
City, Country \\
email address or ORCID}
\and
\IEEEauthorblockN{3\textsuperscript{rd} Given Name Surname}
\IEEEauthorblockA{\textit{dept. name of organization (of Aff.)} \\
\textit{name of organization (of Aff.)}\\
City, Country \\
email address or ORCID}
\and
\IEEEauthorblockN{4\textsuperscript{th} Given Name Surname}
\IEEEauthorblockA{\textit{dept. name of organization (of Aff.)} \\
\textit{name of organization (of Aff.)}\\
City, Country \\
email address or ORCID}
\and
\IEEEauthorblockN{5\textsuperscript{th} Given Name Surname}
\IEEEauthorblockA{\textit{dept. name of organization (of Aff.)} \\
\textit{name of organization (of Aff.)}\\
City, Country \\
email address or ORCID}
\and
\IEEEauthorblockN{6\textsuperscript{th} Given Name Surname}
\IEEEauthorblockA{\textit{dept. name of organization (of Aff.)} \\
\textit{name of organization (of Aff.)}\\
City, Country \\
email address or ORCID}
}

\maketitle

\begin{abstract}
Marine monitoring records are often sparse in space and time and may be asynchronously missing across variables. We propose a small-sample multivariate reconstruction framework that uses a same-time local Gaussian prior and learns residual corrections from temporal states, spatial context, and cross-variable relations. The method is evaluated on nine marine regions with seven environmental variables using chronological train/validation/test splits and shared masks across all methods. At 70\% random missingness, the proposed model is competitive with local interpolation, EOF-style reconstruction, and a joint convolutional recurrent baseline, while the best method changes with the variable and missing-data geometry. Ablation results identify the local prior as the most consistent information source and show that temporal and cross-variable branches provide variable-dependent corrections. The study therefore treats reconstruction as a conditional field-recovery problem and reports the operating conditions under which each information source is useful.
\end{abstract}

\begin{IEEEkeywords}
Marine environment; spatiotemporal data gaps; multivariate interpolation; DINEOF; optimal interpolation; evaluation protocol; applicable boundary
\end{IEEEkeywords}

\section{Introduction}
Marine equipment is exposed to coupled marine and atmospheric conditions, including temperature, salinity, humidity, atmospheric SO$_2$, dissolved oxygen, and pH. These variables affect thermal loading, corrosion, coating degradation, and the construction of environmental load spectra. Monitoring records, however, are incomplete because of sparse trajectories, sensor outages, communication interruptions, and nonuniform sampling. The resulting task is conditional recovery of a multivariate field from partially observed space--time data, rather than unconstrained numerical filling.

Optimal interpolation and EOF methods remain useful for small samples because they are inexpensive and interpretable~\cite{lorenc1986,beckers2003,alveraazcarate2005,beckers2006}. Their covariance or low-rank assumptions can nevertheless oversmooth local gradients and do not directly represent nonlinear temporal or cross-variable dependencies. Deep models such as DINCAE, GRIN, MTGNN, and PriSTI provide richer representations~\cite{barth2020dincae,cini2022grin,wu2020mtgnn,liu2023pristi}, but their behavior depends strongly on sample size and missing-data geometry.

This paper addresses the resulting small-sample regime with a local-prior residual framework. A same-time Gaussian estimate supplies the observation-supported baseline; a ConvGRU state, spatial context, and learned seven-variable relations then predict a correction to that baseline, while visible measurements are copied unchanged. We evaluate this design under shared masks, strict chronological isolation, and several missingness geometries to identify which information sources are effective and when.

\section{Method}

\subsection{Mechanism}

\begin{figure}[t]
    \centering
    \includegraphics[width=1\linewidth]{figures/mechanism.jpg}
    \caption{Overview of the proposed reconstruction mechanism. Visible observations provide a local spatial prior; temporal, spatial-context, and variable-relation branches predict a residual correction, while observed values are copied unchanged.}
    \label{fig:mechanism}
\end{figure}

As shown in Fig.~\ref{fig:mechanism}, visible observations first define a same-time local Gaussian baseline. Temporal state, spatial context, and learned seven-variable relations then predict a residual correction, while observed entries are copied unchanged. The following equations define these operations.

\begin{enumerate}
    \item \textbf{Sparse seven-variable field and mask.}
    The model receives a sparse seven-variable marine environmental field
    together with its visibility mask:
    \begin{equation}
    \mathbf{X}_t \in \mathbb{R}^{H\times W\times 7},\qquad
    \mathbf{M}_t \in \{0,1\}^{H\times W\times 7}.
    \label{eq:input}
    \end{equation}
    The seven variables are sea-surface temperature, air temperature,
    salinity, relative humidity, atmospheric $\mathrm{SO_2}$ concentration,
    dissolved oxygen concentration, and pH.

    \item \textbf{Local prior.}
    A spatial prior field is constructed from visible observations in the
    neighborhood of the same time step using a local Gaussian kernel:
    \begin{equation}
    \mathbf{P}_t(i,j)
    =
    \frac{
    \displaystyle\sum_{(u,v)\in\mathcal{N}(i,j)}
    K_{\sigma}(i-u,j-v)\,
    \mathbf{M}_t(u,v)\mathbf{X}_t(u,v)
    }{
    \displaystyle\sum_{(u,v)\in\mathcal{N}(i,j)}
    K_{\sigma}(i-u,j-v)\,
    \mathbf{M}_t(u,v)+\varepsilon
    },
    \label{eq:local-prior}
    \end{equation}
    where $\mathcal{N}(i,j)$ denotes the local neighborhood of position
    $(i,j)$, $K_{\sigma}$ is the Gaussian kernel, and $\varepsilon$ is a
    small stabilizing constant.

    \item \textbf{Temporal state.}
    Historical feature states and the current temporal phase encoding are
    used to model the temporal evolution of the environmental field:
    \begin{equation}
\mathbf{H}_t
    =
    \operatorname{ConvGRU}
    \left(
    \mathbf{F}_t,\mathbf{H}_{t-1},
    \boldsymbol{\phi}_t
    \right),
\label{eq:temporal-state}
\end{equation}
    where $\mathbf{F}_t$ denotes the current local features,
    $\mathbf{H}_{t-1}$ is the historical hidden state, and
    $\boldsymbol{\phi}_t$ is the temporal phase encoding.

    \item \textbf{Spatial context.}
    Geographic coordinates and local observation density are incorporated as
    spatial contextual information:
    \begin{equation}
\mathbf{S}_t(i,j)
    =
\mathbf{W}_2\,\sigma\!\left(\mathbf{W}_1[\lambda_i,\varphi_j,\rho_t(i,j)]^{\mathsf T}+\mathbf{b}_1\right)+\mathbf{b}_2,
\label{eq:spatial-context}
\end{equation}
    where $\lambda_i$ and $\varphi_j$ denote longitude and latitude,
    respectively, and $\rho_t(i,j)$ represents the local observation density
    around position $(i,j)$.

    \item \textbf{Learned variable relation.}
    A $1\times1$ convolution, or an equivalent channel-mixing layer, jointly
    learns the coupling relationships among the seven environmental
    variables:
    \begin{equation}
\mathbf{V}_t
    =
    \operatorname{Conv}_{1\times1}
    \left(
    \mathbf{X}_t,\mathbf{M}_t
    \right).
\label{eq:variable-relation}
\end{equation}
    This branch represents statistical associations among temperature,
    salinity, humidity, dissolved oxygen, pH, and the other environmental
    variables.

    \item \textbf{Residual fusion.}
    Local spatial priors, temporal states, spatial context, and learned
    variable relations are fused to predict a residual correction to the
    local prior:
    \begin{equation}
\mathbf{R}_t
    =
    f_{\theta}
    \left(
    \mathbf{P}_t,\mathbf{H}_t,\mathbf{S}_t,\mathbf{V}_t
    \right),
\label{eq:residual}
\end{equation}
    \begin{equation}
\widehat{\mathbf{X}}_t
    =
    \mathbf{P}_t+\mathbf{R}_t.
\label{eq:prediction}
\end{equation}

    \item \textbf{Observed-value copy.}
    Original visible observations are preserved without modification, while
    only missing locations are replaced by model predictions:
    \begin{equation}
\widetilde{\mathbf{X}}_t
    =
    \mathbf{M}_t\odot\mathbf{X}_t
    +
    \left(1-\mathbf{M}_t\right)\odot
    \widehat{\mathbf{X}}_t,
\label{eq:observed-copy}
\end{equation}
    where $\odot$ denotes element-wise multiplication.

    \item \textbf{Completed field and diagnostics.}
    The model produces a completed seven-variable environmental field
    $\widetilde{\mathbf{X}}_t$. The reconstruction is subsequently assessed
    using posterior diagnostics, including residual distributions,
    inter-variable correlations, spatial continuity, temporal continuity,
    and physical plausibility.
\end{enumerate}

\section{Experiments}

\subsection{Dataset and Evaluation Protocol}

The dataset contains nine marine regions, with 24 temporal fields collected for
each region. Each field is represented on a common geographic grid and includes
seven environmental variables: sea-surface temperature (\texttt{tos}, K), air
temperature (\texttt{tas}, K), salinity (\texttt{sos}, psu), relative humidity
(\texttt{hurs}, \%), atmospheric $\mathrm{SO_2}$ concentration (\texttt{so2},
dataset units), dissolved oxygen concentration (\texttt{o2}, mg~L$^{-1}$),
and pH (\texttt{ph}, dimensionless). The data
were chronologically divided into 14 training fields, 5 validation fields, and
5 test fields for each region. No field from the validation or test periods was
used for parameter fitting, normalization, or model selection.

To evaluate reconstruction under realistic monitoring conditions, missing
values were generated by masking observed entries. The primary experiment used
a 70\% random missing rate. Additional random-missing experiments were
performed with missing rates from 10\% to 90\%. For each missing rate, five
shared mask seeds (20260906--20260910) were used for all methods, ensuring that
the compared models received exactly the same visible observations and were
evaluated on the same hidden entries. Missingness was sampled independently
for different variables, so that the missing pattern of one environmental
variable did not determine the pattern of another variable.

Four structured missing scenarios were additionally considered: continuous
temporal gaps, grid-boundary missingness, asynchronous variable missingness, and
missingness concentrated in the top 20\% high-gradient areas. The grid-boundary
scenario refers to the boundary of the gridded computational domain and is not
interpreted as a physical coastline. These scenarios were used as stress tests
for different observation geometries rather than as additional training data.

\subsection{Training, Reconstruction, and Model Selection}

For each training field, an internal self-supervised holdout was generated from
the available observations. The model received the remaining visible values,
the corresponding visibility mask, geographic coordinates, observation-density
features, and temporal phase information. The loss was computed only on the
artificially hidden target entries:
\begin{equation}
\mathcal{L}_{\mathrm{train}}
=
\frac{1}{|\Omega_{\mathrm{tar}}|}
\sum_{(i,j,k)\in\Omega_{\mathrm{tar}}}
\left(
\widehat{X}_{t,i,j,k}-X_{t,i,j,k}
\right)^2,
\label{eq:training-loss}
\end{equation}
where $\Omega_{\mathrm{tar}}$ denotes the target entries hidden for
self-supervised training. Consequently, visible values were not treated as
prediction targets, and the model could not use hidden target values as input.

The same training procedure was applied independently to each method. Model
parameters were optimized using AdamW with a learning rate of
$2\times 10^{-3}$ and a weight decay of $10^{-5}$. The gradient norm was clipped
to 1.0. Training was run for at most 55 epochs, with early stopping when the
validation loss did not improve for 10 consecutive epochs. The model checkpoint
with the lowest validation loss was retained for testing. All trainable models
were implemented in PyTorch 2.11.0 and executed with CUDA 12.8 on an NVIDIA
GeForce RTX 5060 Laptop GPU.

The mean and standard deviation used for normalization were estimated only from
valid values in the 14 training fields. The same transformation was then
applied to the validation and test fields, and all reported errors were
converted back to the original physical units. During testing, model parameters
were frozen. For each test field, only the observations that remained visible
under the test mask were provided to the model. The hidden test values were
never used for input construction, model selection, or parameter updates and
were used only for final evaluation. The final completed field was obtained by
copying visible observations directly and replacing only missing entries:
\begin{equation}
\widetilde{\mathbf{X}}_t
=
\mathbf{M}_t\odot\mathbf{X}_t
+
\left(1-\mathbf{M}_t\right)\odot
\widehat{\mathbf{X}}_t.
\label{eq:observed-copy-test}
\end{equation}

\subsection{Compared Methods}

The proposed method, denoted as \textsc{Ours}, was compared with four reference
approaches under the same data split, input masks, target masks, and evaluation
procedures:

\begin{itemize}
    \item \textbf{Climatology}: reconstruction from the training-period
    climatological mean at each grid location and variable;
    \item \textbf{OI analogue}: local observation-based objective-analysis
    reconstruction using the nearest available observations;
    \item \textbf{EOF-style}: low-rank empirical orthogonal function
    reconstruction fitted using the incomplete training fields;
    \item \textbf{JointConvGRU}: a joint multi-variable convolutional recurrent
    network without the proposed local-prior residual formulation;
    \item \textbf{\textsc{Ours}}: the proposed model combining a local Gaussian
    spatial prior, convolutional feature extraction, ConvGRU temporal states,
    temporal phase encoding, coordinate and observation-density features,
    learned seven-variable relations, and residual fusion.
\end{itemize}

All learning-based methods were trained using the same chronological split,
self-supervised target construction, validation-based checkpoint selection,
normalization procedure, and optimization budget. Non-learning reference
methods were fitted using training-period data only and applied directly to the
same masked validation and test fields.

% Figure mapping in D:/research/会议2:
% 图片2.png: random-missingness sensitivity
% 图片3.png: missingness-geometry comparison
% 图片4.png: OURS ablation
% 图片5.png: random-70% main comparison

\subsection{Illustrative Multivariable Reconstructions}
Fig.~\ref{fig:raw-mask} shows the original and independently masked fields for
one test case; Fig.~\ref{fig:imputation} shows the corresponding \textsc{Ours}
reconstruction with shared color limits within each variable.

\begin{figure}[t]
    \centering
    \includegraphics[width=\linewidth]{figures/raw_mask.png}
    \caption{Example of the seven-variable reconstruction input. The first row
    contains the original fields and the second row contains the independently
    masked observations supplied to the model. Each column corresponds to one
    environmental variable.}
    \label{fig:raw-mask}
\end{figure}

\begin{figure}[t]
    \centering
    \includegraphics[width=\linewidth]{figures/imputation.png}
    \caption{Example reconstruction for the same test case. Completed fields are
    shown for all seven variables using the proposed method. Visible input values
    are copied unchanged, and only masked entries are replaced.}
    \label{fig:imputation}
\end{figure}

\section{Results}

\subsection{Main Results under 70\% Random Missingness}

Figure~\ref{fig:random70} presents the region-level RMSE under the primary
70\% random-missingness condition. Errors are reported in the original physical
units for each variable, and the error bars represent 95\% intervals calculated
with marine regions as the clustering units. Because the seven variables have
different physical units and numerical scales, their RMSE values are compared
only within each variable.

OURS achieved the lowest regional mean RMSE for four of the seven variables:
0.4271~K for air temperature (\texttt{tas}), 0.1954~psu for salinity
(\texttt{sos}), 2.0745\% for relative humidity (\texttt{hurs}), and 0.4695 for
the atmospheric $\mathrm{SO_2}$ indicator. JointConvGRU obtained the lowest
RMSE for sea-surface temperature (\texttt{tos}), with an RMSE of 0.2676~K,
whereas the RMSE of OURS was 0.2682~K. OI analogue performed best for dissolved
oxygen (\texttt{o2}) and pH, with RMSE values of 0.0392~mg~L$^{-1}$ and
0.0051, respectively. The differences between OURS and the two strongest
reference methods were generally small for \texttt{tas}, \texttt{sos},
\texttt{so2}, and \texttt{tos}.

Confidence intervals overlap for most variables, indicating variable-specific
rather than universal gains.

\begin{figure}[t]
    \centering
    \includegraphics[width=\linewidth]{figures/random70.png}
    \caption{Region-level RMSE under 70\% random missingness. Each panel
    corresponds to one environmental variable. Gray points denote individual
    marine-region results, colored points denote regional means, and vertical
    bars indicate 95\% intervals.}
    \label{fig:random70}
\end{figure}

\subsection{Sensitivity to the Random Missingness Rate}

Figure~\ref{fig:missingness-sweep} shows RMSE as the random missingness rate
increases from 10\% to 90\%. Error generally grows with missingness, with
larger changes for \texttt{hurs}, \texttt{tas}, and \texttt{tos}.

OURS remained the lowest-error method for \texttt{hurs} across the evaluated
missingness rates. It was also competitive for \texttt{sos}, \texttt{tas}, and
\texttt{so2}, particularly at low and intermediate missingness rates. At 90\%
missingness, JointConvGRU achieved lower RMSE for \texttt{o2} and
\texttt{sos}, EOF-style achieved the lowest RMSE for \texttt{tas}, and OI
analogue was best for \texttt{so2} and pH. JointConvGRU generally performed
best for \texttt{tos} up to 70\% missingness, while OURS became comparable or
slightly better at 90\%. These crossings indicate that the relative ranking
changes as the amount of same-time spatial support decreases.

\begin{figure}[t]
    \centering
    \includegraphics[width=\linewidth]{figures/missingness_sweep.png}
    \caption{Sensitivity of reconstruction RMSE to the random missingness rate.
    Results are shown separately for the seven environmental variables in their
    original units.}
    \label{fig:missingness-sweep}
\end{figure}

\subsection{Performance under Structured Missingness}

Figure~\ref{fig:geometry} summarizes the lowest-error method for each variable
under four structured missingness conditions: continuous temporal gaps, grid
domain-edge missingness, asynchronous variable missingness, and missingness
concentrated in the top 20\% highest-gradient cells. The figure reports
categorical winners for each variable.

The winning method changes across temporal gaps, domain edges, asynchronous
variables, and high-gradient cells. OI and EOF remain strong for broad or
temporally supported fields, while OURS is competitive for selected variables.

\begin{figure}[t]
    \centering
    \includegraphics[width=\linewidth]{figures/geometry.jpg}
    \caption{Lowest-error method for each environmental variable under five
    missingness geometries. The random-70\% column corresponds to the primary
    experiment; the remaining columns represent structured stress tests.}
    \label{fig:geometry}
\end{figure}

\subsection{Component Ablation}

Figure~\ref{fig:ablation} reports the percentage change in RMSE after removing
one component from OURS under a fixed training and masking seed.

Removing the local-neighborhood branch increased the RMSE for \texttt{sos},
\texttt{tas}, \texttt{so2}, and \texttt{o2} by approximately 5.4\%, 3.1\%,
6.8\%, and 5.9\%, respectively. This branch therefore provided the most
consistent contribution in the tested configuration. The effects of removing
the variable-relation, coordinate, temporal-phase, and observation-density
branches depended on the target variable. For example, removing the
coordinate branch reduced the pH RMSE by 6.1\% but increased the \texttt{hurs}
RMSE by 3.6\%; removing the density branch increased the pH RMSE by 5.9\% but
reduced the \texttt{hurs} RMSE by 2.6\%. These mixed directions indicate that
the additional branches provide variable-dependent corrections rather than
uniform gains.

OURS contains 108,176 trainable parameters versus 106,327 for JointConvGRU.

\begin{figure}[t]
    \centering
    \includegraphics[width=\linewidth]{figures/ablation.png}
    \caption{Component ablation of OURS relative to the complete model. Values
    show the percentage change in RMSE after removing one component. Positive
    values indicate increased error.}
    \label{fig:ablation}
\end{figure}

\section{Conclusion}
This paper presented a small-sample framework for reconstructing multivariate marine environmental fields from partially observed space--time records. A same-time local Gaussian prior supplies the observation-supported baseline, while temporal state, spatial context, and learned seven-variable relations predict residual corrections; visible measurements are copied unchanged. Under shared random and structured masks, the best method depended on the variable and missing-data geometry. The local prior was the most consistent component in ablation, whereas temporal and cross-variable branches provided variable-dependent supplements. These results support conditional field recovery as a practical formulation for sparse environmental monitoring. Future work should evaluate larger observation archives and incorporate external circulation and coastline information.

\begin{thebibliography}{00}
\bibitem{lorenc1986} A. C. Lorenc, ``Analysis methods for numerical weather
prediction,'' \emph{Quarterly Journal of the Royal Meteorological Society},
vol. 112, pp. 1177--1194, 1986.
\bibitem{beckers2003} J.-M. Beckers and M. Rixen, ``EOF calculations and data
filling from incomplete oceanographic data sets,'' \emph{Journal of Atmospheric
and Oceanic Technology}, vol. 20, pp. 1839--1856, 2003,
doi: 10.1175/1520-0426(2003)020<1839:ECADFF>2.0.CO;2.
\bibitem{alveraazcarate2005} A. Alvera-Azc{\'a}rate, A. Barth, M. Rixen, and
J.-M. Beckers, ``Reconstruction of incomplete oceanographic data sets using
empirical orthogonal functions: application to the Adriatic Sea surface
temperature,'' \emph{Ocean Modelling}, vol. 9, pp. 325--346, 2005,
doi: 10.1016/j.ocemod.2004.08.001.
\bibitem{beckers2006} J.-M. Beckers, A. Barth, and A. Alvera-Azc{\'a}rate,
``DINEOF reconstruction of clouded images including error maps: application to
the sea-surface temperature around Corsican Island,'' \emph{Ocean Science},
vol. 2, pp. 183--199, 2006, doi: 10.5194/os-2-183-2006.
\bibitem{barth2020dincae} A. Barth \emph{et al.}, ``DINCAE 1.0: a convolutional
neural network with error estimates to reconstruct sea surface temperature
satellite observations,'' \emph{Geoscientific Model Development}, vol. 13,
pp. 1609--1622, 2020, doi: 10.5194/gmd-13-1609-2020.
\bibitem{cini2022grin} A. Cini, I. Marisca, and C. Alippi, ``Filling the G\_ap\_s:
multivariate time series imputation by graph neural networks,'' in
\emph{Proc. International Conference on Learning Representations}, 2022.
\bibitem{wu2020mtgnn} Z. Wu \emph{et al.}, ``Connecting the dots: multivariate
time series forecasting with graph neural networks,'' in \emph{Proc. ACM
SIGKDD}, 2020, pp. 753--763, doi: 10.1145/3394486.3403118.
\bibitem{liu2023pristi} M. Liu \emph{et al.}, ``PriSTI: a conditional diffusion
framework for spatiotemporal imputation,'' in \emph{Proc. IEEE International
Conference on Data Engineering}, 2023, pp. 1927--1939,
doi: 10.1109/ICDE55515.2023.00150.
\end{thebibliography}

\end{document}
